\documentclass[UTF8,a4paper,11pt]{ctexart}

\usepackage{amsmath,amssymb,mathtools}
\usepackage{geometry}
\usepackage{fancyhdr}
\usepackage{enumitem}
\usepackage{hyperref}

\geometry{left=2.5cm,right=2.5cm,top=2.5cm,bottom=2.5cm}
\setlength{\parindent}{2em}
\linespread{1.2}
\setlength{\headheight}{13.6pt}

\pagestyle{fancy}
\fancyhf{}
\fancyfoot[C]{\thepage}
\fancyhead[L]{数值分析\quad 第八次理论作业\quad 学号：241098117，姓名：张城宗}
\fancyhead[R]{习题四、五}

\newcommand{\dd}{\,\mathrm{d}}
\newcommand{\ee}{\mathrm e}
\newcommand{\QED}{\hfill$\square$}

\begin{document}

\begin{center}
  {\LARGE\bfseries 数值分析\quad 第八次理论作业}\\[6pt]
  {\large 习题四第 21、28、34、35、37、38 题；习题五第 3、4 题}\\[4pt]
  {\normalsize 学号：241098117\quad 姓名：张城宗}
\end{center}

\bigskip

%=====================================================
\noindent\textbf{习题四\quad 第 21 题}\quad
设函数 $f(x)$ 在 $[-1,1]$ 上有六阶连续导数，$p_5(x)\in \mathbb R_5[x]$ 是满足
\[
  p_5(x_i)=f(x_i),\qquad p_5'(x_i)=f'(x_i),
  \qquad x_i=-1,0,1
\]
的 Hermite 插值多项式。试证明
\[
\begin{aligned}
  \int_{-1}^{1}f(x)\dd x
  \approx \int_{-1}^{1}p_5(x)\dd x
  &=
  \frac{7}{15}f(-1)+\frac{16}{15}f(0)+\frac{7}{15}f(1)  \\
  &\quad+\frac{1}{15}f'(-1)-\frac{1}{15}f'(1),
\end{aligned}
\]
并推导其余项。

\medskip
\noindent\textbf{证明}\quad
由于节点和积分区间关于原点对称，设
\[
  \int_{-1}^{1}p_5(x)\dd x
  =A[f(-1)+f(1)]+Bf(0)+C f'(-1)-C f'(1).
\]
该公式应对所有不超过五次的多项式精确成立。分别令
$f(x)=1,x^2,x^4$，得
\[
  2A+B=2,\qquad
  2A-4C=\frac23,\qquad
  2A-8C=\frac25.
\]
解得
\[
  A=\frac{7}{15},\qquad B=\frac{16}{15},\qquad C=\frac{1}{15}.
\]
而奇次多项式的积分为零，上式由对称性也自动给出零值，因此对
$0,1,\ldots,5$ 次多项式均精确成立。于是对 Hermite 插值多项式
$p_5$ 有
\[
  \int_{-1}^{1}p_5(x)\dd x
  =
  \frac{7}{15}f(-1)+\frac{16}{15}f(0)+\frac{7}{15}f(1)
  +\frac{1}{15}f'(-1)-\frac{1}{15}f'(1).
\]

由 Hermite 插值余项，
\[
  f(x)-p_5(x)
  =
  \frac{f^{(6)}(\xi_x)}{6!}
  (x+1)^2x^2(x-1)^2,\qquad -1<\xi_x<1.
\]
因为 $(x+1)^2x^2(x-1)^2\ge 0$，由积分中值定理，存在
$\xi\in(-1,1)$，使得
\[
\begin{aligned}
  R(f)
  &=\int_{-1}^{1}\bigl(f(x)-p_5(x)\bigr)\dd x  \\
  &=\frac{f^{(6)}(\xi)}{720}
    \int_{-1}^{1}x^2(x^2-1)^2\dd x  \\
  &=\frac{f^{(6)}(\xi)}{720}
    \left(\frac{2}{7}-\frac{4}{5}+\frac{2}{3}\right)
   =\frac{1}{4725}f^{(6)}(\xi).
\end{aligned}
\]
所以
\[
  \boxed{
  \int_{-1}^{1}f(x)\dd x
  =
  \frac{7}{15}f(-1)+\frac{16}{15}f(0)+\frac{7}{15}f(1)
  +\frac{1}{15}f'(-1)-\frac{1}{15}f'(1)
  +\frac{1}{4725}f^{(6)}(\xi).}
\]
\QED

\bigskip\bigskip

%=====================================================
\noindent\textbf{习题四\quad 第 28 题}\quad
应用自适应 Simpson 求积法计算积分
\[
  I=\int_{0}^{3}x\sqrt{1+x^2}\dd x,
\]
要求误差不超过 $10^{-2}$。

\medskip
\noindent\textbf{解}\quad
记 $f(x)=x\sqrt{1+x^2}$。在区间 $[a,b]$ 上的 Simpson 值记为
\[
  S(a,b)=\frac{b-a}{6}
  \left[f(a)+4f\!\left(\frac{a+b}{2}\right)+f(b)\right].
\]
先在 $[0,3]$ 上计算
\[
  S(0,3)=10.1517434034.
\]
将区间二分后得到
\[
  S(0,1.5)+S(1.5,3)=10.2012724869.
\]
自适应 Simpson 的误差估计为
\[
  \frac{|S(0,1.5)+S(1.5,3)-S(0,3)|}{15}
  =0.0033019389<10^{-2}.
\]
因此不必继续细分。若采用 Richardson 修正，则
\[
\begin{aligned}
  I
  &\approx S(0,1.5)+S(1.5,3)
  +\frac{S(0,1.5)+S(1.5,3)-S(0,3)}{15} \\
  &=10.2045744258.
\end{aligned}
\]
所以
\[
  \boxed{I\approx 10.20457.}
\]
作为校验，精确值为
\[
  \int_0^3x\sqrt{1+x^2}\dd x
  =\frac{(1+x^2)^{3/2}}{3}\bigg|_0^3
  =\frac{10\sqrt{10}-1}{3}\approx 10.20759,
\]
实际误差约为 $3.02\times 10^{-3}$，满足题目要求。

\bigskip\bigskip

%=====================================================
\noindent\textbf{习题四\quad 第 34 题}\quad
给出计算积分
\[
  \int_{-1}^{1}f(x)(1+x^2)\dd x
\]
的两点插值求积公式，使它的代数精确度为 $3$，并导出求积公式的离散误差。

\medskip
\noindent\textbf{解}\quad
这是权函数 $\rho(x)=1+x^2$ 下的两点 Gauss 型求积公式。设二次正交多项式为
\[
  \omega_2(x)=x^2-\alpha.
\]
由正交条件
\[
  \int_{-1}^{1}(x^2-\alpha)(1+x^2)\dd x=0
\]
可得
\[
  \alpha=
  \frac{\int_{-1}^{1}x^2(1+x^2)\dd x}
       {\int_{-1}^{1}(1+x^2)\dd x}
  =
  \frac{\frac23+\frac25}{2+\frac23}
  =\frac25.
\]
因此求积节点为
\[
  x_1=-\sqrt{\frac25},\qquad x_2=\sqrt{\frac25}.
\]
由对称性 $A_1=A_2=A$，再令 $f(x)=1$，得
\[
  2A=\int_{-1}^{1}(1+x^2)\dd x=\frac83,
  \qquad A=\frac43.
\]
所以所求公式为
\[
  \boxed{
  \int_{-1}^{1}f(x)(1+x^2)\dd x
  \approx
  \frac43 f\!\left(-\sqrt{\frac25}\right)
  +\frac43 f\!\left(\sqrt{\frac25}\right).}
\]

下面求离散误差。设
\[
  E(f)=\int_{-1}^{1}f(x)(1+x^2)\dd x
  -\frac43 f\!\left(-\sqrt{\frac25}\right)
  -\frac43 f\!\left(\sqrt{\frac25}\right).
\]
由两点 Gauss 型求积公式的余项形式，
\[
  E(f)=\frac{f^{(4)}(\xi)}{4!}
  \int_{-1}^{1}(1+x^2)
  \left(x^2-\frac25\right)^2\dd x,\qquad -1<\xi<1.
\]
而
\[
  \int_{-1}^{1}(1+x^2)
  \left(x^2-\frac25\right)^2\dd x
  =\frac{136}{525}.
\]
故
\[
  \boxed{E(f)=\frac{17}{1575}f^{(4)}(\xi),\qquad -1<\xi<1.}
\]

\bigskip\bigskip

%=====================================================
\noindent\textbf{习题四\quad 第 35 题}\quad
求 Gauss 型求积公式
\[
  \int_0^1 f(x)\ln x\dd x
  \approx A_1f(x_1)+A_2f(x_2)
\]
的系数 $A_1,A_2$ 以及节点 $x_1,x_2$，并导出离散误差。

\medskip
\noindent\textbf{解}\quad
令
\[
  \omega_2(x)=x^2+ax+b.
\]
它应关于权函数 $\ln x$ 与 $1,x$ 正交。由
\[
  \int_0^1 x^k\ln x\dd x=-\frac{1}{(k+1)^2}
\]
得
\[
  \begin{cases}
  \displaystyle
  \int_0^1 (x^2+ax+b)\ln x\dd x=0,\\[2mm]
  \displaystyle
  \int_0^1 x(x^2+ax+b)\ln x\dd x=0.
  \end{cases}
\]
即
\[
  \frac19+\frac{a}{4}+b=0,\qquad
  \frac1{16}+\frac{a}{9}+\frac{b}{4}=0.
\]
解得
\[
  a=-\frac57,\qquad b=\frac{17}{252}.
\]
因此节点为方程
\[
  x^2-\frac57x+\frac{17}{252}=0
\]
的两个根：
\[
  \boxed{
  x_1=\frac{15-\sqrt{106}}{42},\qquad
  x_2=\frac{15+\sqrt{106}}{42}.}
\]

再由公式对 $1,x$ 精确成立，
\[
  A_1+A_2=-1,\qquad
  A_1x_1+A_2x_2=-\frac14,
\]
解得
\[
  \boxed{
  A_1=-\frac{2\sqrt{106}+9}{4\sqrt{106}},
  \qquad
  A_2=\frac{9-2\sqrt{106}}{4\sqrt{106}}.}
\]
于是
\[
  \boxed{
  \int_0^1 f(x)\ln x\dd x
  \approx
  -\frac{2\sqrt{106}+9}{4\sqrt{106}}
  f\!\left(\frac{15-\sqrt{106}}{42}\right)
  +\frac{9-2\sqrt{106}}{4\sqrt{106}}
  f\!\left(\frac{15+\sqrt{106}}{42}\right).}
\]

设离散误差为 $E(f)=I(f)-I_1(f)$。由两点 Gauss 型公式余项，
\[
  E(f)
  =
  \frac{f^{(4)}(\xi)}{4!}
  \int_0^1
  \left(x^2-\frac57x+\frac{17}{252}\right)^2\ln x\dd x,
  \qquad 0<\xi<1.
\]
直接计算得
\[
  \int_0^1
  \left(x^2-\frac57x+\frac{17}{252}\right)^2\ln x\dd x
  =-\frac{647}{226800}.
\]
所以
\[
  \boxed{E(f)=-\frac{647}{5443200}f^{(4)}(\xi),\qquad 0<\xi<1.}
\]

\bigskip\bigskip

%=====================================================
\noindent\textbf{习题四\quad 第 37 题}\quad
应用三点 Gauss--Legendre 求积公式计算积分
\[
  \int_0^1\frac{\sin x}{1+x}\dd x.
\]

\medskip
\noindent\textbf{解}\quad
三点 Gauss--Legendre 公式在 $[-1,1]$ 上为
\[
  \int_{-1}^{1}F(t)\dd t
  \approx
  \frac59F\!\left(-\sqrt{\frac35}\right)
  +\frac89F(0)
  +\frac59F\!\left(\sqrt{\frac35}\right).
\]
作变换
\[
  x=\frac{t+1}{2},\qquad \dd x=\frac12\dd t.
\]
记 $r=\sqrt{3/5}$，则
\[
\begin{aligned}
  \int_0^1\frac{\sin x}{1+x}\dd x
  &\approx
  \frac{5}{18}
  \frac{\sin\frac{1-r}{2}}{1+\frac{1-r}{2}}
  +\frac49\frac{\sin\frac12}{1+\frac12}
  +\frac{5}{18}
  \frac{\sin\frac{1+r}{2}}{1+\frac{1+r}{2}}  \\
  &\approx 0.2842484986.
\end{aligned}
\]
故
\[
  \boxed{\displaystyle
  \int_0^1\frac{\sin x}{1+x}\dd x\approx 0.28424850.}
\]

\bigskip\bigskip

%=====================================================
\noindent\textbf{习题四\quad 第 38 题}\quad
作适当变换，应用 Gauss--Chebyshev 求积公式计算积分
\[
  I=\int_1^3 x\sqrt{4x-x^2-3}\dd x
\]
的准确值。

\medskip
\noindent\textbf{解}\quad
注意到
\[
  4x-x^2-3=1-(x-2)^2.
\]
令
\[
  t=x-2,\qquad x=t+2,\qquad \dd x=\dd t,
\]
则
\[
  I=\int_{-1}^{1}(t+2)\sqrt{1-t^2}\dd t
  =
  \int_{-1}^{1}
  \frac{(t+2)(1-t^2)}{\sqrt{1-t^2}}\dd t.
\]
因此可以使用 Gauss--Chebyshev 公式，其中
\[
  F(t)=(t+2)(1-t^2)
\]
是三次多项式。$n$ 点 Gauss--Chebyshev 公式的代数精确度为
$2n-1$，故只要 $2n-1\ge3$，即 $n\ge2$，就能得到准确值。

取 $n=2$，节点为
\[
  t_1=\frac{1}{\sqrt2},\qquad t_2=-\frac{1}{\sqrt2},
\]
权系数均为 $\pi/2$。于是
\[
\begin{aligned}
  I
  &=\frac{\pi}{2}
  \left[
  F\!\left(\frac1{\sqrt2}\right)
  +F\!\left(-\frac1{\sqrt2}\right)
  \right]  \\
  &=\frac{\pi}{2}
  \left[
  \left(2+\frac1{\sqrt2}\right)\frac12
  +\left(2-\frac1{\sqrt2}\right)\frac12
  \right]
  =\pi.
\end{aligned}
\]
所以
\[
  \boxed{I=\pi.}
\]

\bigskip\bigskip

%=====================================================
\noindent\textbf{习题五\quad 第 3 题}\quad
假设函数 $g(x)$ 与 $h(x)$ 在区间 $[a,b]$ 上连续，证明初值问题
\[
  y'=g(x)y+h(x),\qquad y(a)=\eta
\]
在 $[a,b]$ 上有唯一解，并且对任何初始值都是适定的。

\medskip
\noindent\textbf{证明}\quad
令
\[
  f(x,y)=g(x)y+h(x).
\]
因为 $g,h$ 在 $[a,b]$ 上连续，所以 $f$ 在带形区域
$D=\{(x,y):a\le x\le b,\ y\in\mathbb R\}$ 上连续。又
\[
  |f(x,y_1)-f(x,y_2)|
  =|g(x)|\,|y_1-y_2|
  \le L|y_1-y_2|,
\]
其中
\[
  L=\max_{a\le x\le b}|g(x)|.
\]
所以 $f$ 关于 $y$ 满足 Lipschitz 条件。由初值问题解的存在唯一性定理，
对任意初值 $\eta$，该初值问题在 $[a,b]$ 上存在唯一解。

也可以写出显式解。令
\[
  G(x)=\int_a^x g(s)\dd s,
\]
则
\[
  y(x)=\ee^{G(x)}
  \left[
  \eta+\int_a^x \ee^{-G(t)}h(t)\dd t
  \right].
\]
若初值由 $\eta$ 改为 $\tilde\eta$，对应解为 $y,\tilde y$，则
\[
  y(x)-\tilde y(x)=\ee^{G(x)}(\eta-\tilde\eta),
\]
从而
\[
  |y(x)-\tilde y(x)|
  \le \ee^{L(b-a)}|\eta-\tilde\eta|.
\]
这说明解连续依赖于初始值。若右端项中的 $h$ 也有微小扰动，则由同样的积分表达式可得
\[
  \|y-\tilde y\|_\infty
  \le
  \ee^{L(b-a)}
  \left(
  |\eta-\tilde\eta|+(b-a)\|h-\tilde h\|_\infty
  \right),
\]
故解也连续依赖于右端项。因此该初值问题对任意初始值都是适定的。
\QED

\bigskip\bigskip

%=====================================================
\noindent\textbf{习题五\quad 第 4 题}\quad
试用 Taylor 级数法（取 $p=3$）导出求解初值问题
\[
  y'=\frac{1}{1+y^2},\qquad y(0)=1
\]
的数值方法。

\medskip
\noindent\textbf{解}\quad
设 $x_n=nh$，$y_n\approx y(x_n)$。三阶 Taylor 方法为
\[
  y_{n+1}
  =
  y_n+h y'_n+\frac{h^2}{2}y''_n+\frac{h^3}{6}y'''_n.
\]
由微分方程
\[
  y'=\frac{1}{1+y^2}
\]
可得
\[
  y''
  =
  \frac{\dd}{\dd x}\left(1+y^2\right)^{-1}
  =
  -\frac{2y}{(1+y^2)^3}.
\]
继续求导，
\[
\begin{aligned}
  y'''
  &=
  \frac{\dd}{\dd x}
  \left[-2y(1+y^2)^{-3}\right]  \\
  &=
  \left[-2(1+y^2)^{-3}+12y^2(1+y^2)^{-4}\right]y'  \\
  &=
  \frac{10y^2-2}{(1+y^2)^5}.
\end{aligned}
\]
代入三阶 Taylor 展开，得到数值格式
\[
  \boxed{
  y_{n+1}
  =
  y_n
  +\frac{h}{1+y_n^2}
  -\frac{h^2y_n}{(1+y_n^2)^3}
  +\frac{h^3(5y_n^2-1)}{3(1+y_n^2)^5},
  \qquad y_0=1.}
\]
这就是取 $p=3$ 时的 Taylor 级数法。

\end{document}
